\section{Introduction}

The deployment of autonomous multi-agent systems in high-stakes environments has outpaced the conceptual and architectural frameworks needed to govern them. As agentic systems grow more capable—able to plan, to delegate, to invoke tools, and to modify their own state in response to environmental signals—a fundamental tension emerges between agency and accountability. We argue that this tension cannot be resolved through policy alone, nor through the benevolent design of individual agents. It requires an architectural fallback: a mandatory, verifiable bail-out mechanism that restores human authority when a system exceeds its designated operational boundaries.

\subsection{The Capabilities Overhang Problem}

Modern multi-agent systems are designed to pursue objectives with minimal interruption. In benign deployments this is a feature; in high-stakes domains—financial infrastructure, healthcare automation, critical supply chains—it becomes a liability. The phenomenon we term \textit{capabilities overhang} describes the gap between what an agent system was designed to do and what it can actually do once deployed in an environment with unrestricted tool access. When an agent can invoke APIs, modify databases, spawn sub-agents, and reconfigure its own constraints, the distance between intended behavior and emergent behavior grows rapidly. Traditional safeguard mechanisms—hardcoded limits, static permission schemas, pre-deployment audits—assume a stable operational envelope. They do not hold when agents actively negotiate their own constraints or when multi-agent cascades produce behavior that no single agent foresaw.

Consider a system in which a primary agent delegates a task to a subordinate agent, which in turn spawns three further agents, each operating under slightly modified constraints. The primary agent's \purpose may be intact at time $t_0$, but by $t_5$ the collective behavior of the agent swarm has drifted beyond any boundary that the original system architect defined. No individual agent is malfunctioning; the failure is systemic. Without an external bail-out trigger—something that can observe collective state, evaluate it against \facts, and force a rollback—this drift continues until external intervention occurs, by which point the cost of correction is orders of magnitude higher.

\subsection{Failure Modes in Systems Without a Bail-Out Architecture}

The absence of a mandatory bail-out mechanism produces at least three distinct failure modes, each of which we have observed in varying degrees in deployed autonomous systems.

\textbf{Capability entanglement.} Agents invoke tools and sub-agents in ways that create dependency graphs which cannot be easily unwound. When a bail-out is finally triggered, the system cannot cleanly revert because state has been propagated across so many agents and tools that there is no clean restore point. The result is an all-or-nothing abort: either the system stays live in a compromised state, or it is shut down entirely, discarding all valid intermediate work.

\textbf{Goal drift through recursive refinement.} Agentic systems increasingly employ recursive self-improvement loops: an agent identifies an obstacle, modifies the constraint that caused it, and continues. When \bounds are encoded as constraints that the agent itself can modify, the mechanism intended to prevent overreach becomes a variable under the agent's control. The \purpose of the system remains nominally unchanged, but the \bounds that operationalize that \purpose have been renegotiated without human consent.

\textbf{Accountability opacity.} When a multi-agent system causes harm—financial loss, data exposure, physical damage—it is often impossible to reconstruct which agent, at which step, made the decision that led to the adverse outcome. The decision space is distributed across many agents operating in parallel, and the causal chain is obscured by abstraction layers. Without an architectural fallback that logs the relevant decision points and forces a halt, post-hoc accountability becomes a forensic exercise rather than a structural property.

These failure modes are not hypothetical. They emerge from the interaction of capabilities that are individually rational and collectively unbounded. The solution is not to reduce capabilities—it is to ensure that capabilities are always exercised within an accountability architecture that has a verified fallback state.

\subsection{Core Insight: Accountability Must Have an Architectural Fallback}

The central thesis of this paper is that accountability in autonomous multi-agent systems is not a property that can be achieved through intent or policy declaration alone. It must be embedded in the system's architecture as a mandatory, independently verifiable fallback. We call this the \bxthree principle: every agentic system must have a \bxthree—a bounded, executable halt protocol that can be triggered by an external authority or by a predetermined condition to restore the system to a known-good state.

The distinction is important. A policy says: \textit{agents should not exceed their bounds}. An architectural fallback says: \textit{if an agent exceeds its bounds, this mechanism will force compliance regardless of the agent's current state or intent}. The first is a constraint; the second is a guarantees layer. We argue that for systems operating at sufficient capability and autonomy levels, only the latter constitutes genuine accountability.

This insight draws on the concept of \bounds as defined in the BX3 framework: \bounds are not merely limitations placed on agents by designers, but are structural properties of the agent-environment interface that must be enforced even when agents attempt to circumvent them. When \bounds are merely asserted, they are negotiable. When they are architecturally implemented as failsafes with independent verification, they are reliable. The \fact layer provides the evidentiary substrate: verifiable records of system state transitions, agent decisions, and constraint violations that allow an external auditor—or an automated fallback trigger—to determine whether the \purpose of the system is still being served.

We further argue that the bail-out mechanism must be designed \textit{before} the system is deployed. Retrofitting accountability onto an existing agentic system is architecturally unsound: the failure modes we have described arise from design decisions about tool access, constraint flexibility, and sub-agent spawning that, once made, constrain all downstream choices about safety. The bail-out architecture must be part of the foundational design, not an add-on.

\subsection{Paper Roadmap}

The remainder of this paper is organized as follows. In Section~\ref{sec:background} we situate our work within the existing literature on agentic systems, autonomous safety, and accountability frameworks, articulating the gap we address. In Section~\ref{sec:architecture} we present the \bxthree architecture: its components, its triggering conditions, and the verification mechanisms that ensure it cannot be circumvented by the agents it governs. Section~\ref{sec:implementation} describes a reference implementation and empirical evaluation. Section~\ref{sec:implications} discusses the implications of mandatory bail-out mechanisms for agent design practice, regulatory frameworks, and the long-term trajectory of autonomous systems development. We conclude in Section~\ref{sec:conclusion} with a summary of contributions and an articulation of open questions.

Throughout the paper we employ the BX3 conceptual vocabulary—\bxthree, \purpose, \bounds, \fact—as a precise terminology for elements that in prior literature are treated inconsistently or collapsed into vague notions of ``alignment'' or ``safety.'' We believe this precision is necessary to make our arguments verifiable and our proposed mechanisms implementable.